\chapter{Tonsillectomy}
\section*{Overview}
Tonsillectomy removes the palatine tonsils from the back of the throat.
\section*{Common Instruments and Equipment}
The exact setup varies by technique and setting. Common items may include:
\begin{itemize}[nosep,leftmargin=*]
  \item mouth gag and surgical light
  \item tonsil dissection instruments or powered device
  \item suction and electrosurgical hemostasis equipment
\end{itemize}

\section*{Why It Is Done}
It may be performed for recurrent severe tonsillitis, obstructive sleep-disordered breathing, peritonsillar abscess, or selected tonsil disease.
\section*{How It Is Performed}
Under general anesthesia, the tonsils are separated from the surrounding tissue and removed through the mouth. Adenoidectomy may be performed at the same operation.
\section*{Contraindications}
No single contraindication applies to every patient; the treating team assesses whether the procedure is suitable for the individual and the reason for treatment.

\section*{Precautions}
The care team reviews relevant health conditions, medicines, allergies, and procedure-specific risks, and explains any changes needed for the individual.

\section*{Risks and Recovery}
Postoperative bleeding, dehydration, pain, infection, nausea, and anesthesia complications are possible. Hydration, pain control, and urgent attention to bleeding are important during recovery.
\section*{References}
\begin{enumerate}\item MedlinePlus. Tonsillectomy and Adenoidectomy. U.S. National Library of Medicine; 2025.\item Current Medical Diagnosis and Treatment. McGraw-Hill; 2025.\end{enumerate}
\clearpage
